\subsection{All odd-place quartic subset characters}

\begin{statement}{A.118}{Proved}
{Quartic negative-subset evaluations vanish at every odd place}
Let $p$ be an odd prime and let
$x,r,d_1,\ldots,d_k,w_1,\ldots,w_k\in\Q_p^\times$ satisfy
\[
 w_i^4=x+(1-x)d_i\quad(1\le i\le k),
       \qquad \prod_{i=1}^k d_i=-r^2.
\tag{A.118.1}
\]
Then
\[
                         (x,2r)=0
                              \quad\hbox{in }\operatorname{Br}(\Q_p).
\tag{A.118.2}
\]
The same assertion holds over $\R$.
There is no bound on the number $k$ of factors.

More geometrically, fix any finite set of distinct rational
punctures $c_1,\ldots,c_k$ and a quartic cover
\[
       C:\quad T-c_i=a_iX_i^4,\qquad a_i\in\Q^\times.
\tag{A.118.3}
\]
Let $c\in\Q\setminus\{c_1,\ldots,c_k\}$, and let $B_c$ be
the largest abelian field contained in every field factor of
the actual boundary fibre
\[
       \Q[Z_1^{\pm1},\ldots,Z_k^{\pm1}]
             /(Z_i^4-(c-c_i)/a_i).
\tag{A.118.4}
\]
For every character $\chi$ of $\operatorname{Gal}(B_c/\Q)$,
the evaluation of $(T-c,\chi)$ is constant on the ordinary
points of $C(\Q_p)$ with $T\ne c$, for every odd $p$.
It is also constant on the corresponding real points.

In particular, for the rationally pointed cover
$a_i=R-c_i$, where $R\in\Q\setminus\{c_i,c\}$, every
normalized evaluation is zero at all these places:
\[
              \left(\frac{T-c}{R-c},\chi\right)=0.
\tag{A.118.5}
\]
For this pointed cover the exact finite-character test of
A.85, at exponent four and with any finite set of prescribed
ordinary old points, therefore reduces to its dyadic
evaluations.  This last reduction does not assert that the
dyadic evaluations vanish.
\end{statement}
\begin{proof}
If $x=1$ or if $x$ is a square, (A.118.2) is immediate
from the quadratic norm description of a quaternion.
We henceforth assume that $x$ is a nonsquare.

\emph{Primes at which a fourth root of unity is present.}
Suppose first that $p\equiv1\pmod4$, and fix $i\in\Q_p$
with $i^2=-1$.
Put $K=\Q_p$, choose $\alpha$ with $\alpha^4=x$, and set
$L=K(\alpha)$.
This is a splitting field, since its roots are
$\alpha,i\alpha,-\alpha,-i\alpha$.
Its Galois group embeds in the cyclic group of order four
by $\sigma\mapsto\sigma(\alpha)/\alpha$.
A proper subgroup would fix $\alpha^2$, making $x$ a square
in $K$.  Thus $L/K$ is cyclic of degree four, with quadratic
subfield $K(\sqrt{x})$.
The literal norm identities are
\[
 d_i=\frac{w_i^4-x}{1-x}
     =N_{L/K}\left(\frac{w_i-\alpha}{1-\alpha}\right),
 \qquad -r^2=(ir)^2\in N_{L/K}(L^\times).
\tag{A.118.6}
\]
By the local reciprocity isomorphism
\cite[Theorem 4.1.2]{BIR_KedlayaCFT}, the image of $ir$ in
$\operatorname{Gal}(L/K)\simeq\Z/4$ has order at most two.
The subgroup of order two is exactly the kernel of the
restriction map to $\operatorname{Gal}(K(\sqrt{x})/K)$.
The same reciprocity map therefore shows that $ir$ is a
norm from $K(\sqrt{x})$.
The explicit identity
\[
                     \frac{2r}{ir}=-2i=(1-i)^2
\tag{A.118.7}
\]
makes $2r$ a norm from this quadratic extension as well.
This proves (A.118.2) in this case, with the actual
cyclic-to-quadratic restriction map retained.

\emph{Primes at which the fourth roots of unity are unramified quadratic.}
Suppose now that $p\equiv3\pmod4$.
If $v_p(x)<0$, replace the data by
\[
 x'=x^{-1},\qquad w_i'=w_i^{-1},\qquad
 d_i'=d_i/w_i^4,\qquad r'=r/\prod_i w_i^2.
\tag{A.118.8}
\]
Direct substitution gives (A.118.1) for the primed data.
The first quaternion entry is inverted and the second is
multiplied by a square, so its Brauer class is unchanged.
We may therefore write $x=p^h u$ with $h\ge0$ and $u$ a
unit.  Write $s_i=4v_p(w_i)\in4\Z$, and let $\lambda$ be
the quadratic character of $\F_p^\times$.

First suppose $h\notin4\Z$.
Then $h>0$, the denominator $1-x$ is a unit with residue
one, and none of the $s_i$ equals $h$.
If $s_i<h$, then $v_p(d_i)=s_i$ and its normalized unit
has square residue.  If $s_i>h$, then $v_p(d_i)=h$ and its
normalized unit has the residue of $-u$.
Let $b$ be the number of indices of the latter type.
The normalized-unit and valuation parts of
$\prod_i d_i=-r^2$ give respectively
\[
       (-\lambda(\bar u))^b=-1,
       \qquad
       2v_p(r)=\sum_{s_i<h}s_i+bh.
\tag{A.118.9}
\]
Here $\lambda(-1)=-1$ has been used explicitly.
The first equality forces $b$ to be odd and
$\lambda(\bar u)=1$.
The second then forces $h$ to be even.
Every principal unit is a square at an odd prime, by
Hensel lifting, so these two conclusions would make $x$
a square.  This contradicts our standing assumption.

Consequently $h\in4\Z$.
Since $x$ remains a nonsquare, $\bar u$ is a nonsquare.
The denominator $1-x$ is a unit: this is automatic for
$h>0$, and for $h=0$ it follows from $\bar u\ne1$.
When $s_i=h$, the two leading units in $w_i^4-x$ cannot
cancel, because one is a square and the other is a
nonsquare.  Hence in every case
\[
             v_p(d_i)=\min(s_i,h)\in4\Z,
                    \qquad v_p(r)\in2\Z.
\tag{A.118.10}
\]
The extension $K(\sqrt{x})/K$ is unramified quadratic.
Its norm group consists exactly of the elements with even
valuation: the norm on residue units is onto, and its
successive lifts on principal units are onto because the
residue-field trace is onto; the norm of the common
uniformizer is its square.
Since $p$ is odd, $v_p(2r)=v_p(r)$ is even.
Thus $2r$ is a norm, proving (A.118.2).

Over $\R$, a negative $x$ in (A.118.1) would force every
$d_i>-x/(1-x)>0$, contrary to their negative product.
Thus $x>0$ and the real quaternion is zero.

\emph{The exact map from two old points to (A.118.1).}
Let $P_R=(R,\xi_1,\ldots,\xi_k)$ and
$P_T=(T,\eta_1,\ldots,\eta_k)$ be two ordinary local
points of (A.118.3), with $R,T\ne c$.
If $R=T$, their character evaluations are identical.
Otherwise put
\[
 x=\frac{T-c}{R-c},\qquad
 w_i=\eta_i/\xi_i,\qquad
 D_i=\frac{c-c_i}{a_i},\qquad
 d_i=\frac{c-c_i}{R-c_i}=D_i/\xi_i^4.
\tag{A.118.11}
\]
There is the exact identity
\[
                 w_i^4=\frac{T-c_i}{R-c_i}
                        =x+(1-x)d_i.
\tag{A.118.12}
\]
For the individual quadratic characters one has $d_i\ne1$
and the actual quadratic norm
\[
 x=\frac{w_i^4-d_i}{1-d_i}
    =N_{K(\sqrt{d_i})/K}
          \left(\frac{w_i^2-\sqrt{d_i}}{1-\sqrt{d_i}}\right)
\tag{A.118.13}
\]
when $d_i$ is a nonsquare; for square $d_i$ its character
is zero.  Thus $(x,D_i)=(x,d_i)=0$ at every place.

By A.96, all remaining generators of the common character
group of (A.118.4) are the characters of $2r_J$, where
\[
 \prod_{i\in J}D_i=-r_J^2,\qquad r_J\in\Q^\times,
 \qquad
 \prod_{i\in J}d_i
      =-\left(r_J/\prod_{i\in J}\xi_i^2\right)^2.
\tag{A.118.14}
\]
The new local root differs from $r_J$ by a square.
Applying (A.118.2) to precisely this subset shows
$(x,2r_J)=0$ at every odd prime and at the real place.
Together with (A.118.13), this proves the claimed local
constancy for every common character, including all
relations among the displayed generators.

For $a_i=R-c_i$, use the rational point with every
$\xi_i=1$ in (A.118.11); this gives (A.118.5).
For any adelic choice of old ordinary points avoiding $c$,
subtract the global constant class $(R-c,\chi)$.
Global reciprocity makes the sum of the constant's local
invariants zero.  Every normalized odd and real invariant
has just been proved zero, so the character sum appearing
in A.85 is exactly its dyadic normalized invariant.
The same statement holds when finitely many of these old
points are prescribed, which is the finite initialization
scope asserted here.
If $c=R$, omitted from this normalization, the boundary
fibre itself has the rational root tuple $(1,\ldots,1)$.
One of its field factors is then $\Q$, so its common
abelian field is $\Q$ and there are no nonzero common
characters to test.
\end{proof}

The certificate
\path{code/odd_quartic_subset_certificate.py}
checks the exact norm and inversion identities and tests
the actual odd Hilbert evaluations over finite collections
of local fourth-power classes and their generated subset
products.  Its retained output records the range and
counts; these finite checks supplement the proof above.
